% TOP_PROBLEM: {"id":30007704,"problem_number":"LOCAL-30007704","title":"Calibrating sensitivity assumptions","category_id":21,"category":{"id":21,"name":"economics","display_name":"Economics","description":"Open economic theory statements, published research agendas, and proposed scoped specifications; question provenance and historical priority are qualified per record.","slug":"economics","order_index":21,"created_at":"2026-10-08T00:00:00Z"},"status":"research_question","external_url":null,"legacy_ids":[],"published":false,"statement":"Calibrate known omitted-variable sensitivity bounds using substantively defensible confounder information and valid external checks.","economics":{"original_id":193,"field":"Econometrics, causal inference and measurement","kind":"identification","formalization_basis":"adapted_research_question","source_status":"research_agenda","priority_status":"earliest_found","priority_note":"Earliest relevant explicit question or agenda passage located in this audit. No claim of first-ever formulation; earlier versions and later solutions were not exhaustively traced.","earliest_verified_reference":{"authors":["Carlos Cinelli","Chad Hazlett"],"title":"Making Sense of Sensitivity: Extending Omitted Variable Bias","year":2020,"url":"https://academic.oup.com/jrsssb/article/82/1/39/7056023","doi":"10.1111/rssb.12348","locator":"Introduction, paragraph describing the third fundamental obstacle to sensitivity analysis","evidence_summary":"Identifies connecting sensitivity calculations to substantive understanding as a fundamental obstacle."},"current_reference":{"authors":["Carlos Cinelli","Chad Hazlett"],"title":"Making Sense of Sensitivity: Extending Omitted Variable Bias","year":2020,"url":"https://academic.oup.com/jrsssb/article/82/1/39/7056023","doi":"10.1111/rssb.12348","locator":"Introduction, paragraph describing the third fundamental obstacle to sensitivity analysis","evidence_summary":"Identifies connecting sensitivity calculations to substantive understanding as a fundamental obstacle."},"formalization_note":"The source verifies a broader question or research agenda; the population, estimand, equations and restrictions here are our scoped adaptation. This is not a quoted canonical conjecture, and the audit does not prove contemporary unsolved status for every special case. Independent math review tightened feasibility, existence, or estimand assumptions where required."},"statusQualification":"Published question or research agenda verified at the cited locator. The mathematical specification is adapted for this catalogue and includes additional assumptions; source publication does not establish first-ever priority or certify this exact model remains unsolved. The source verifies a broader question or research agenda; the population, estimand, equations and restrictions here are our scoped adaptation. This is not a quoted canonical conjecture, and the audit does not prove contemporary unsolved status for every special case. Independent math review tightened feasibility, existence, or estimand assumptions where required.","statusReviewedAt":"2026-10-08"}
% FIRST_POSTED: 2026-10-08
% ENTRY_KIND: statement_only
% !TeX program = lualatex
\documentclass[11pt]{article}
\usepackage{fontspec}
\usepackage[a4paper,margin=25mm]{geometry}
\usepackage{amsmath,amssymb,mathtools}
\usepackage{xurl}
\usepackage[hidelinks]{hyperref}
\newcommand{\sref}[2]{\hyperlink{source:#1:#2}{[S#2]}}
\setlength{\emergencystretch}{3em}
\begin{document}
% problemId:30007704
\section*{Calibrating sensitivity assumptions}\label{30007704}
\subsection{Definitions and mathematical statement}
Consider a study population with outcome \(Y\), treatment \(D\), measured covariates \(X\), and an unmeasured scalar confounder \(U\). Assume a constant-effect linear causal model \(Y=\beta D+g(X)+\gamma U+\varepsilon\), with \(g\) linear in the recorded covariate basis and \(E[\varepsilon\mid D,X,U]=0\); otherwise the coefficient need not equal a population treatment effect. Let \(\beta_{\rm obs}\) be the coefficient from a specified linear projection omitting \(U\). Define \(r_D\) as the partial squared correlation of \(D\) and \(U\) after projecting on \(X\), and \(r_Y\) as that of \(Y\) and \(U\) after projecting on \(D,X\). Assume the partial correlations are well-defined and both displayed residual standard deviations are positive. If \(r_D<1\), the omitted-variable bias satisfies
\[
|\beta_{\rm obs}-\beta|=\frac{\sigma_{Y\mid D,X}}{\sigma_{D\mid X}}\sqrt{\frac{r_Yr_D}{1-r_D}},
\]
where the two \(\sigma\)'s are population residual standard deviations for the stated projections. Known sensitivity algebra therefore supplies bounds once credible bounds on \(r_D,r_Y\) are given. The research task is to calibrate those bounds using external validation samples, measured confounders, negative controls or substantive restrictions, without assuming measured and unmeasured confounders are interchangeable. State validity conditions for each calibration source, propagate estimation uncertainty, and evaluate coverage and informativeness across a declared class of confounding mechanisms. The agenda concerns credible calibration, not rediscovery of the bias identity.

\par\textbf{Formulation and scope.} The source verifies a broader question or research agenda; the population, estimand, equations and restrictions here are our scoped adaptation. This is not a quoted canonical conjecture, and the audit does not prove contemporary unsolved status for every special case. Independent math review tightened feasibility, existence, or estimand assumptions where required.

\par\textbf{Open-question provenance.} An explicit question or research agenda was verified at the source locator. This does not by itself certify that every later version remains unresolved \sref{30007704}{1}.

\par\textbf{Historical priority.} Earliest relevant explicit question or agenda passage located in this audit. No claim of first-ever formulation; earlier versions and later solutions were not exhaustively traced.

\subsection{Short English statement}
Calibrate known omitted-variable sensitivity bounds using substantively defensible confounder information and valid external checks.

\subsection{Sources}
\begin{itemize}\raggedright
\item[\textnormal{[S1]}]\hypertarget{source:30007704:1}{} Carlos Cinelli, Chad Hazlett. 2020. Making Sense of Sensitivity: Extending Omitted Variable Bias. Introduction, paragraph describing the third fundamental obstacle to sensitivity analysis.
\url{https://academic.oup.com/jrsssb/article/82/1/39/7056023}.
DOI: 10.1111/rssb.12348.
{\small\textit{Earliest verified question/agenda reference; historical priority qualified below. Identifies connecting sensitivity calculations to substantive understanding as a fundamental obstacle.}}
\end{itemize}
\end{document}
